\section{Theta congruences at a characteristic-zero point}
\label{sec:theta}

Fix the odd prime $p$ and the fields $K,F,L$ chosen in
Section~\ref{sec:classical}.  This section supplies the analytic input
for the theta implication.  Its input is a zero of the undepleted
series $b_w^0$ of Section~\ref{sec:pair}, at a characteristic-zero point
of the deformation space.  Its output is a sequence of integral
ordinary theta sections whose constant Fourier--Jacobi coefficients
tend uniformly to zero at that point.  The next section supplies a
nonconstant coefficient of bounded valuation and extracts a Selmer
class.  Both parts must use the same finite group characters, including
the characters of increasing conductor at $p$.

We retain the CM type induced by the distinguished place $w$ of $K$.
Write $\mathcal C$ for the algebraic character of the norm-one torus
of $K/\Q$ with infinity type $s\mapsto s$, and also for its norm
pullback to $L/F$.  Its finite conductor is supported at fixed split
places away from $p$.  Its $p$-avatar is denoted by
$\widehat{\mathcal C}$.  In the distinguished unit position coordinates
the avatar contributes the first power of the coordinate; after norm
pullback it contributes the product of the two coordinates.  This is
the character and the reciprocity normalization of
\cite[Section 2, ``Coefficient, character, and measure conventions'']{C}.
Fix a sufficiently narrow weight branch.  We require that difference
powers of $\widehat{\mathcal C}$ on this branch factor through the
free anticyclotomic group $\Gamma$ used in Section~\ref{sec:pair}.
Indeed a fixed power kills its finite conductors away from $p$ and its
torsion character; the anticyclotomic ray-class group unramified away
from $p$ has free rank one.  The resulting character on $\Gamma$ has
nonzero logarithmic direction, as is seen on the distinguished local
units.  Narrowing this branch does not restrict how large its positive
integer weights can be.

\subsection{Auxiliary tests and the point of evaluation}

The compact theta source and the target are the unitary groups of
\[
 A=(L^2,1_2),\qquad
 W=B\oplus\mathbb H,\qquad
 B=(L^2,b_0\sqrt D\,1_2),\qquad
 \mathbb H=Le\oplus Lf.
\]
Here $A$ is Hermitian, $W$ is skew-Hermitian, $b_0\in F^{\times,+}$
is a $p$-unit, and the line coefficients have the even valuation
required for matched self-dual lattices at nonsplit finite places.
The auxiliary-field conditions, including $(D/|D'|)=1$, give this
choice by the norm-invariant calculation in
\cite[Equations (4.1)--(4.2)]{C}.  The target has signature $(3,1)$ at
both real places of $F$.  Use the trace additive character and unit
self-dual position coordinates at $p$.  Order the distinguished
columns as $(e,B_1,B_2,f)$, with $e$ in the multiplicative direction.

Let $\pi$ be the compact automorphic representation obtained from the
weight-two form of $E$ by real quadratic base change and
Jacquet--Langlands transfer.  Its infinity type and central character
are trivial.  Its standard base change to $L$ is cuspidal, because the
auxiliary fields avoid the inducing field in the CM case.  Its
elliptic local parameters are tempered.  At nonsplit finite places it
has the spherical vectors prescribed by the matched lattices; the
level primes split completely in $L$.  These are precisely the
cuspidality and local hypotheses used for the compact theta and
binary period constructions in \cite[Sections 4--5]{C}.  We make no
assumption on $\pi_p$.

\begin{lemma}[Auxiliary tests]\label{theta:auxiliary-data}
One may choose two sets of fixed finite theta data with the following
properties.  Their extra rational good primes requiring nonspherical
target controls are split completely in $L$, are different from $p$,
and the two sets are disjoint.  Outside these sets and the elliptic
bad primes, the target tests are spherical.  At an unramified-twist
Steinberg place the source vector is Iwahori invariant and the target
test has line-stabilizer parahoric level.  At an additive place fixed
nonspherical tests are permitted.

For either set there are a compact vector $\varphi\in\pi^\vee$ and a
finite-order character $d_0$ of
$[T]=T(F)\backslash T(\mathbb A_F)$, where $T=\mathrm U(Le_2)$,
such that
\[
 I_T(\varphi,d_0)=\int_{[T]}\varphi(s)d_0(s)\,ds\ne0.
\]
They retain the stated spherical and Iwahori conditions, and $d_0$ is
unramified there.  For the two lines $LB_j$ in $B$, let
$\vartheta_{j,d}$ denote the unary theta function with normalized
position monomial of degree $d$.  Its splitting character contributes
a fixed shift $c_j$ to the infinity weight, with $c_1+c_2=1$.
There are compatible characters $\xi_{j,d}$ of weight $d+c_j$
such that, when $d_1+d_2+1=m$,
\[
 \xi_{1,d_1}\xi_{2,d_2}=\mathcal C^m d_0,\qquad
 T_j(d)=\int_{[\mathrm U(1)]}
          \vartheta_{j,d}(s)\xi_{j,d}(s)^{-1}\,ds.
\]
The tests with unit position support at $p$ can be chosen so that
$T_1(d_{10})$ and $T_2(d_{20})$ are separately nonzero at positive
integers $d_{10},d_{20}$.  The integer $d_{10}+d_{20}+1$ is divisible by the
fixed torsion orders and by an arbitrarily prescribed fixed power
of $p$.
\end{lemma}

\begin{proof}
The ray-character construction allows its additional conductors at
arbitrarily chosen sufficiently deep split places: imposing that the
ratio of a principal idele to its conjugate be $1$ at such a place
removes the finite unit obstruction, after which the character extends
over the finite ray-class quotient.  The same construction gives the
weight and inverse unary splitting characters.  Neat level may also
be imposed at a new split place.  Carry out these choices twice,
avoiding the finite set already used.

The required spherical tests at an unmarked split place can be seen
in matrix coordinates.  Choose integral $e,B_1,B_2$ columns, and
scale the $f$ column dually to the additive conductor.  Their mixed
position lattice contains the primitive norm-one orbit.  Conjugate
the hyperspecial to these bases if necessary.  On a split unary
line, the integral lattice test is torus-unit invariant and its
diagonal Mellin integral is the nonzero central Tate integral.
The absolute-value half shift in the oscillator and temperedness
give its nonzero value; fixed volume and uniformizer powers are
nonzero.  Thus these local tests may be retained in the unary
nonvanishing construction.  At nonsplit places use the matched
hyperspecial tests.  Refining a metaplectic multiplier or cusp level
for an auxiliary calculation does not impose a new target Hecke
control at that place.

Start with a nonzero compact vector having the required invariances.
Translate at a free split place until its restriction to $T$ is
nonzero.  More explicitly, the torus supplies determinants, and
strong approximation for $\mathrm{SU}(A)$ away from that split place
places the other components in their prescribed compact opens.
Fourier expansion of this nonzero function on the compact torus
quotient gives $d_0$.  Its unramifiedness at the retained compact
places follows from invariance.  A fixed diagonal translation at
$p$ makes the vector invariant under the integral lower unipotent
without making its torus period zero.

Apply the simultaneous unary nonvanishing lemma of
\cite[Lemma 6.2]{C}.
Its inputs here are coherent matched line pairs, unramified data at
nonsplit finite places, self-dual critical unary characters of
parallel infinity weights $d_j+c_j$ with nonnegative remaining type,
a reduction prime greater than two split completely in $L$, and a
split degree-one variation prime.  Those primes can avoid all the
fixed sets just prescribed.  Inverse variation of the two characters
preserves their product.  The split-line calculation above retains
the prescribed unramified tests; the one-dimensional local Hom space
identifies them with the nonzero local functional.  At $p$ use the
unit-character and unit-cut tests of that lemma.  Its congruence
choice gives the asserted sum of weights.  The analyticity branches
depend on powers of $\mathcal C$ and on fixed CM scales, not on the
new finite-order coefficients; hence they can be narrowed, including
the $\Gamma$ condition above, before making the nonvanishing choices.
\end{proof}

Include the fixed split primes in both sets in the exponent conditions
for the moving-prime tuple.  Distinct such primes have linearly
independent exponent vectors, as specified in the tuple construction.
They need not be inserted into the depleted set $S^\circ$: at a
primitive strict test, adding an unramified good place does not change
the complex.

Write $R=\Z_p[[t,\mathbf u]]$ for the deformation ring, with the fixed
coefficient extension understood when needed.  An interior point is
a tuple $z=(t,\mathbf u)$ in the maximal ideal of the integers of a
finite extension of $\Q_p$.  A horizontal divisor is a height-one
prime not containing $p$.

For either elliptic representation, a primitive strict test at a
specified side $a\in\{w,\bar w\}$ uses the strict condition at
$a$, the full condition at the other $p$-place, unramified
conditions at the fixed places away from $p$, and the permitted
full conditions at moving places.  We will use open sets where
the local invariant groups vanish, so the global degree-one
cohomology is the corresponding Selmer group.

The fixed-place exceptional polynomials below have two concrete
sources.  A spanning denominator is a nonzero minor of a finite
matrix expressing specified Schwartz vectors as combinations of
controlled translates.  A binary separation polynomial is a
product of resultants between the two-dimensional elliptic factor
and the final factors $|\cdot|^{-1/2},|\cdot|^{1/2}$ in the standard
theta parameter, after norm-power shifts with exponents between
$-4$ and $4$; its nonvanishing excludes linkage between these
factors by the norm powers occurring in the local separation test.
At a moving
prime the inverse fake Euler factor is the degree-two reciprocal
of the central local Euler factor formed from the bare elliptic
Satake parameters and the scalar uniformizer value, discarding
scalar inertia.  These definitions identify the tests needed for
point selection before their use in the binary comparison.

\begin{lemma}[Choice of a point]\label{theta:test-point}
Let $\mathfrak p$ be a horizontal divisor on which $b_w^0$ vanishes.
Fix a strict side for each of the two elliptic representations, and
suppose both corresponding primitive strict cohomology groups
vanish generically on $\mathfrak p$.  There is an interior point
$z_0=(t_0,\mathbf u_0)$ on this divisor, over a fixed finite extension
$k/\Q_p$, retaining their vanishing and the following properties.

The local $H^0,H^2$ vanishings at $p$ and the acyclicity of the moving
local complexes hold for both elliptic representations, with the
conjugate and $L$-restricted tests as required.  One of the two sets
of auxiliary tests in Lemma~\ref{theta:auxiliary-data} has nonzero
fixed-place spanning and binary separation denominators at $z_0$.
At every moving prime the inverse fake Euler factor can be chosen
nonzero in the limiting test.  Binary separation there can be tested
on inertia when the tested character is ramified, and on Frobenius
when it is unramified.
\end{lemma}

\begin{proof}
The cohomological conditions are open conditions on the divisor:
they are the nonvanishing of minors of the finite free models.
Their generic validity was proved for the horizontal tests in
Sections~\ref{sec:models} and~\ref{sec:pair}.  The asserted absence of
the two jumps is another such open condition.

At a fixed split good prime the additional exceptional conditions
are zeros of finitely many Laurent polynomials in its unramified
scalar parameter.  These are the spanning denominators of
\cite[Lemma 6.4]{C} and the
Frobenius separation resultants of
\cite[Equation (8.11) and Lemma 8.6]{C}.
They are nonzero at trivial scalar and weight zero.  The parameters
at the two $F$-places are powers or inverse powers of the single
parameter at the rational prime.  Two different rational primes
cannot both give exceptional identities on $\mathfrak p$: their
independent exponent vectors make equations fixing both parameters
to algebraic values have codimension at least two.  One may verify
this after a finite flat power change of multiplicative coordinates,
which diagonalizes the two exponent vectors.  At most one rational
prime is therefore identically exceptional, and one of the two
disjoint auxiliary sets avoids it.  Avoid the remaining finitely
many zero loci at the point.

For a moving prime, the inverse fake Euler factor is a polynomial of
degree at most two in the character value of a uniformizer.  Changing
the uniformizer by a unit leaves the cleared formula unchanged:
the only term using the fake factor as an actual Euler factor is
multiplied by the unit norm sum, on which the character is trivial.
If the tested unit character is nontrivial, it has at least three
values, since $p$ is odd.  A unit change avoids the two possible
roots.  If it is identically trivial on the divisor, its own
coordinate is $u_j=0$, and the trace condition in the moving-prime
construction makes the limiting Euler determinant nonzero.
The same trace condition gives the Frobenius separation resultants
in this case.  In the ramified case choose a unit with nontrivial
tested value; it separates the elliptic head from the two trivial
inertia eigenvalues.  Preserve these finitely many nonvanishings.

The interior algebraic points are dense in each nonempty affinoid
open subset of a characteristic-zero divisor.  The bounded-series
specialization argument therefore supplies $z_0$ outside the listed
proper zero loci.  Enlarge its finite field of definition for the
fixed data only.
\end{proof}

Fix such a point.  Label finite ring-class characters so that the
toric integral with weight $\eta^{-1}$ has the Galois character
weight $\chi$ in \eqref{eq:source-10}.  A class character trivial on
the base ideles descends along $z\mapsto z/z^c$ to the norm-one
torus; pulling back that descended character gives the original
class character, rather than its square.  Pull it to $L$ by norm.
There are two simultaneous inverse conventions.  Choose one once
at all places.  In the later Galois extraction it gives
$V\otimes\chi^{\pm1}$; if the inverse occurs, conjugation identifies
its strict condition with the opposite strict condition for $\chi$.
For the construction using $\bar w$ make the conjugate choice.  Thus
the two constructions exchange sides and use the same global scalar
character in the determinant problems.

Let $G_i$ be the product of the finite tame groups and
$\Gamma/\Gamma^{p^{a_i}}$, with $a_i\to\infty$, and let $\nu_i$ be
its universal character.  Denote the group coefficient ring by
$\mathcal O[G_i]$, allowing finite extensions for the values of
Schwartz tests.  Substitution at $z_0$ is a ring homomorphism modulo
a precision tending to infinity.  For example,
$(1+t_0)^{p^{a_i}}\to1$, and the analogous relations hold for the
tame variables.  Integer lifts of exponents give the same value
modulo that precision.  A uniformly bounded group polynomial stays
bounded after substitution; products and limits of bounded series
are respected because tails tend uniformly to zero on every closed
interior polydisc.  Coefficient fields of classical theta sections
may vary with $i$.  Only the field $k$ of the comparison values is
fixed.

\subsection{The ordinary branch and increasing depth}

Set $\eta=\mathcal C^m\nu_i$, where $m>0$ is on the fixed branch.
With the even splitting characters trivial, the upper theta integral is
\[
 \Theta_{i,m}(g)
 =\int_{[\mathrm U(A)]}\theta_{\Phi_{i,m}}(g,h)
       \varphi(h)\eta^{-1}(\det h)\,dh.
\]
The archimedean Fock vector has determinant degree $m-1$.  Its
source type is $\det^m$, including the vacuum determinant.  Its
target weight is $(m,m,1;-1)$.  Use the lattice dual to the monomial
lattice in the $2\times2$ minors, so extraction of a monomial does
not introduce a binomial coefficient.  The normalized creation
operators and CM periods are those of
\cite[Section 4, ``Frames, periods, and unary moments'']{C}: creation
has leading linear position coordinate, and fixed algebraic unit
scales are restored after computing in Igusa differential frames.

At both distinguished $p$-places put
\begin{equation}
 \Phi_p(Y)=\phi_p(Y_{[e,B_1]})
      \nu_{i,p}(\det Y_{[e,B_1]})
      \mathbf 1_{M_{2,2}(\Z_p)}(Y_{[B_2,f]}),
 \qquad \operatorname{supp}\phi_p\subset\mathrm{GL}_2(\Z_p).
 \tag{18}\label{eq:source-18}
\end{equation}
The bare function $\phi_p$ has fixed depth.  The determinant factor
cancels the determinant character in source integration; hence the
source compact and its volume remain fixed.  Choose $s_i$ large
enough for the locally constant tests, and use the subgroup $J(s_i)$
of upper block type $(2,1,1)$: its upper radical and final two unit
groups are full, its head Levi is principal congruence at depth
$s_i$, and its lower blocks have sufficient depth.  This level marks
a rank-two flag and frame inside the rank-three multiplicative
subgroup.  It does not mark a splitting of the etale quotient.

\begin{lemma}[Canonical levels]\label{theta:canonical-levels}
At each finite depth $s$, the ordinary branch of the normalization
in level $J(s)$ on which the rank-three flag is multiplicative adds
only finite etale character-group flags and frames.  At a toroidal
cusp it does not ramify the center Fourier parameter.  Every
component of this branch meets a compatible ordinary cusp, and
reduction of all Fourier--Jacobi expansions detects divisibility of
an integral section by the uniformizer.  These assertions hold
after every finite coefficient extension, with no loss of $p$-adic
precision depending on $s$.
\end{lemma}

\begin{proof}
The multiplicative subgroup of an ordinary $p$-divisible group is
canonical at every finite level.  Cartier duality identifies its
subflags and frames with quotients and markings of a finite etale
character group.  Once the rank-three member is this subgroup, the
remaining level data are exactly these etale data.  The full upper
radical is what removes any choice of a lift of the etale quotient.
The corresponding generic-fiber choices form an open and closed
branch after normalization.  Idempotents extend over the ordinary
completion; normalization commutes with that completion by
excellence.

At a rank-one boundary the semiabelian object is an extension of a
binary abelian object by a torus.  The preimage of its binary
multiplicative torsion is of multiplicative type.  Its character
group and flags are constructed before division by the Raynaud
period lattice, so they do not involve the center Fourier
parameter.  The resulting chart covers are finite etale covers of
the abelian extension-data charts.  Pointing such a cover over a
strict henselian boundary base makes it an abelian scheme and the
covering an isogeny, by lifting addition through the etale cover.
This is the geometric construction in
\cite[Lemma 7.1]{C}, applied
separately at depth $s$; no numerical bound in its proof depends on
$s$.

The hyperspecial base has ordinary cusps on every component.
Each connected finite etale component maps onto its base component
and consequently meets its cusp charts.  The expansion principle
on these charts detects a zero reduced section.  Repeated division
by a uniformizer proves saturation at every precision.  These are
statements about reduction and etale covers, so finite ramification
of the coefficient field changes neither the assertion nor the
valuation normalized by $v_p(p)=1$.
\end{proof}

\begin{lemma}[The two $p$-operators]\label{theta:p-operators}
For every $i,m$, the unnormalized right-coset operators
$U_j=[J(s_i)\operatorname{diag}(1_j,p^{-1}1_{4-j})J(s_i)]$,
$j=2,3$, act on \eqref{eq:source-18} by $p^2$.  In weight
$(m,m,1;-1)$ the integral ordinary normalization is $p^{-2}U_j$,
so both normalized comparison values are $1$.
\end{lemma}

\begin{proof}
The upper shifts are indexed by $M_{j,4-j}(\F_p)$.  On the support
of the test the head has rank two modulo $p$, so the expanding
columns can become divisible by $p$ in exactly
$p^{(j-2)(4-j)}$ ways.  None of these equations changes the head
function, including its determinant character.  The split Weil
action supplies $p^{4-j}$.  Their product is
$p^{(j-1)(4-j)}=p^2$ for $j=2,3$.  The computation uses the full
upper radical and is independent of its opposite depth and of
the head character's conductor.
\end{proof}

At a moving $q$ use the integral open-pivot test
\[
 \mathbf1_{\mathrm{GL}_2(\mathcal O_q)}(Y_{[e,B_1]})
 \eta_q(\det Y_{[e,B_1]})
 \mathbf1_{M_{2,2}(\mathcal O_q)}(Y_{[B_2,f]}).
\]
Its equivariance also cancels the source determinant character on
the full source compact.  In particular source integration never
divides by the order of a growing character quotient.  The fixed
tests elsewhere are those already specified, or one of finitely many
fixed variants.  The spherical parameter at an unmarked split place is
\[
 (\pi\otimes\eta)\oplus|\cdot|^{-1/2}\oplus|\cdot|^{1/2},
\]
by the matched unramified theta calculation of
\cite[Lemma 4.2 and Equation (4.10)]{C}.

\subsection{A moment bound independent of the position cut}

The increasing conductor in \eqref{eq:source-18} requires more than
the fixed-level statement of the unary-moment lemma in
\cite[Section 4]{C}.  We now prove the extension that will be used.
An integral position lattice in this statement means the lattice in
the fixed Igusa frame, before restoring algebraic CM unit scales.
If those scales are instead included in the coordinates, the norm
is taken on the resulting scaled compact lattice; a coefficient-field
unit is not identified with an element of $\Z_p$.

\begin{lemma}[Bounded cuts of unary moments]\label{theta:cut-moments}
Fix the integral position lattices, additive characters, Igusa
frames, and the finite tests away from $p$ used for a finite product
of unary theta functions.  Moving tame tests may have bounded
group-ring coefficients and additive Fourier normalizations that
are $p$-units.  At $p$ let $c$ be any locally constant function on
the position lattices, with
coefficients in any finite extension or finite group algebra and
$\|c\|\leq1$.  Its conductor is unrestricted.

For every finite polynomial combination $P$ of normalized position
moments, the value of the corresponding period-normalized theta
jets at a fixed integral ordinary CM object satisfies
\[
       \bigl|\mathcal J(cP)\bigr|_p
       \leq C\,\|P\|_{\mathrm{pos}}.
\]
The constant $C$ depends only on the fixed data.  It is independent
of the moment degrees, the conductor of $c$, the moving tame levels,
and the coefficient extension.  The same bound applies to the
canonical cusp expansions used to test ordinary integrality.
The uncut moment functional is a bounded measure on the compact
position lattice.  When $c$ is invariant under the coordinate
signs, inserting $c$ multiplies that measure by $c$.
\end{lemma}

\begin{proof}
We first prove a scalar estimate.  If a cut or test has values in
a group algebra, write it in the group basis and apply the scalar
argument to each coefficient.  Coefficients of products are sums
of products, so the ultrametric inequality retains the same bound,
independently of the group order.  In particular we take squares
of scalar coefficient jets, rather than infer coefficientwise
integrality from a square in the group algebra.

Fix a finite list of scalar cuts and polynomials.  All their
conductors and degrees are finite.  Choose the initial integral
position frame with first vector in the multiplicative direction
at the fixed CM object, and use this same finite marking in the
level construction.  For these unary theta
functions use the ordinary Hilbert modular Igusa tower with the
fixed tame Schwartz level and its multiplicative differential
trivializations, as in \cite[Lemma 4.4]{C}.  At each split $p$-factor
the multiplicative height is one.  Marking its character group at
a deeper level is finite etale, and does not mark a lift of the
etale quotient.  Over a toroidal cusp it leaves the Fourier
parameter unramified.  To see which cusps detect this branch,
include the toroidal boundary in this construction.  At a Tate
cusp its multiplicative group is the Tate torus, and its
character-group markings are finite etale unit markings of that
torus.  A connected component of this finite etale cover maps onto
its connected base component and hence has a point over the chosen
ordinary Tate cusp.  Choose its adelic cusp representative integral
at $p$, using Iwasawa decomposition and approximation in the cusp
parabolic, and realizing this level marking.  Relative to the
fixed initial position frame, its marked connected vector is
$ue$ modulo $p^s$ with $u$ a unit, at each Hilbert $p$-factor.
Its integral symplectic comparison matrix $g$ therefore has
lower-left entry divisible by $p^s$.  In a two-dimensional factor,
write $g=\left(\begin{smallmatrix}a&b\\c&d\end{smallmatrix}\right)$.
Since $a,d$ are units and $c\in p^s\Z_p$, explicitly
\[
 g=
 \begin{pmatrix}d^{-1}&b\\0&d\end{pmatrix}
 \begin{pmatrix}1&0\\c/d&1\end{pmatrix}.
\]
The second matrix is congruent to the identity in the lower
direction.  Take $s$ at least the stabilizer depth of the finite
list of tests.  The second matrix then fixes those tests.
Thus every component has a detecting cusp representative in
$B(\Z_p)\operatorname{Stab}(c)$ relative to the original position
frame.  No preliminary symplectic transformation of the cut is
being made.  The Fourier interchange of the two
directions would mark the etale direction instead and does not
describe this canonical branch.  The scalar Hilbert
expansion map on this canonical branch is therefore saturated,
by the same reduction argument as in
Lemma~\ref{theta:canonical-levels}.  This is the scalar Hilbert
tower, rather than the unitary target of that lemma.
We may compare the finite list on one sufficiently deep such
level.  We shall remove every denominator introduced in making
that comparison, rather than bound it in terms of the level.

For integral positions the upper integral unipotent acts through
the additive quadratic phase and is in the stabilizer with our
unit self-dual normalization.  The diagonal stabilizer can be made
as deep as the local constancy of the cuts requires.  Smoothness
supplies a lower stabilizer of some finite depth; equivalently one
may choose that depth after bounding the Fourier-transform support
of this individual finite list.  These statements put all the tests
on the canonical multiplicative levels just described.

Here is the effect on moments of these compatible frame changes.
The preceding construction compares each detecting cusp directly
with the fixed CM position frame through
$B(\Z_p)\operatorname{Stab}(c)$.  Their effective changes are
therefore upper triangular: they
rescale positions by units, permute the tests, and multiply by unit
phases.  For a diagonal unit $u$, the test reads $\phi(ux)$, while
the multiplicative marking changes the differential by $u^{-1}$.
The normalized degree-$d$ moment therefore reads $(ux)^d$ in the
same position coordinate.  An upper unipotent contributes a phase
common to the polynomial moments.  Thus both the cut and the
polynomial are transformed in the same coordinate.  We use no
uniform bound on the Fourier transform of fine cuts at $p$; a
Fourier interchange with the opposite position direction is not
among these compatible changes.

Consequently the coefficient at index $\beta$ of a polynomial
combination at such a cusp is
\begin{equation}\label{theta:position-coefficient}
       \sum_{Q(x)=\beta}\phi(x)c(x)P(x).
\end{equation}
The sum is finite, and $Q$ has unit quadratic coefficient in the
normalized position frame.  It has absolute value at most the
fixed bound for $\phi$ times $\|P\|_{\mathrm{pos}}$.  The same
calculation holds for several positions.  Additive transforms away
from $p$ have normalization a residue-characteristic power or a root
of unity and cost only the fixed denominators already included in
$C$.  Moving determinant characters have unit group coefficients.
The bound in \eqref{theta:position-coefficient} is thus independent
of every changing depth.

We justify that saturation applies to these analytic jets.  At any
individual set of degrees and levels, enlarge the tame multiplier
level, including the level at two, so that products of two base
half-weight theta functions are algebraic integral-weight forms.
Leibniz and Gauss--Manin express products of their normalized jets
as rational expressions in jets of the base products.  There are
no horizontal poles: locally in characteristic zero the square of
a base theta has even divisor, and after removing that divisor an
etale square-root frame makes the connection and its fixed iterates
regular.  Cross-products fix the signs and descend these products.
The logarithmic cusp calculation is the same.  Projection to the
ordinary unit-root splitting therefore makes the square of every
finite linear combination an element of the ordinary tower ring
with $p$ inverted.  At an ordinary CM object this splitting agrees
with the CM splitting.  These algebraicity and comparison facts are
the finite-level facts proved in the unary-moment argument of
\cite[Lemma 4.4]{C}; they require
finiteness of the chosen list, not a uniform bound for its depth.

Let $G$ be such a square, multiplied by the fixed clearing factor
predicted by \eqref{theta:position-coefficient}.  Its expansions
are integral at all compatible cusps.  If $\varpi^rG$ is integral
with $r$ minimal and positive, every expansion of $\varpi^rG$
vanishes modulo $\varpi$.  By the scalar Hilbert expansion
principle just established, its reduction is zero on every
component.  Division by $\varpi$
contradicts minimality.  Thus $r=0$.  Evaluation and taking the
square root give the claimed bound for the original combination.
This argument removes its whole initial vertical denominator, even
when that denominator depends on the cut.  It is unchanged over a
ramified coefficient extension with the valuation normalized by
$v_p(p)=1$.

For a general cut, separate the parity pieces of the finite
tests and polynomial jets with
$2^{-r}\sum_{\varepsilon\in\{\pm1\}^r}
\varepsilon^e[\varepsilon]$.  These projectors have norm at most
one because $p$ is odd.  In each coordinate write the degree as
$e_0+2j$, with $e_0\in\{0,1\}$; the corresponding base unary
theta has weight $1/2+e_0$, and its $j$th normalized jet has
position factor $x^{e_0}Q(x)^j$.  Carry out the preceding square
argument for these parity bases; their cross-products supply
the common descent.  Each expansion still has the form
\eqref{theta:position-coefficient}.  Thus the estimate applies to
arbitrary cuts, including coset cuts obtained by decomposing a
matrix determinant character; it does not require invariance under
the separate signs of its entries.

With the cut equal to one, the estimate gives a bounded linear
functional on all polynomial parity pieces, and polynomial density
extends it to a bounded measure.  For the additional identification
of a cut with multiplication of this fixed measure, assume that
$c$ is sign invariant.  Approximate it uniformly by sign-invariant
polynomials, using symmetrization.  Apply the same bound to the
difference between a cut moment and its polynomial approximation.
In \eqref{theta:position-coefficient} that difference is exactly
$\phi(x)(c(x)-P_n(x))P(x)$ in the same frame.  Its norm tends to
zero.  Hence the measure with the cut is $c$ times the uncut
measure.  All these arguments are coefficientwise for finite group
algebras, so their bounds do not involve the group order.
\end{proof}

Positive degrees tending to infinity remove the nonunit positions.
On the units, fixing the torsion component makes
$x^d=\omega(x)^d\exp(d\log\langle x\rangle)$ analytic in $d$ on a
sufficiently narrow branch.  Lemma~\ref{theta:cut-moments} makes this
convergence uniform for bounded cuts.  Restoring fixed CM unit-scale
powers preserves the bound and is analytic after narrowing the
branch once.  This is the interpolation of unary periods used for
the nonconstant coefficient in the next section.

\begin{proposition}[Integral theta sections]\label{theta:integral-family}
There is a fixed integer $c\geq0$ such that $p^c\Theta_{i,m}$ is an
integral ordinary section on the canonical branch, coefficientwise
in $\mathcal O[G_i]$.  The same $c$ works for all increasing depths
$s_i$, all sufficiently large positive weights on the fixed branch,
all the fixed test variants, and bounded locally constant cuts in
the head of \eqref{eq:source-18} that retain the source determinant
equivariance on a fixed compact.  The two trailing columns keep
their full integral lattice tests; the canonical cusp changes
preserve their integral-translation invariance, as proved below.
Its
Fourier--Jacobi coefficients are integral in the monomial-dual
differential lattice.
\end{proposition}

\begin{proof}
Choose canonical cusp representatives integral at $p$ and preserving
the multiplicative member modulo the chosen deep level.  In the
mixed expansion the $f$ column is Fourier-dualized to the position
$x$.  It remains invariant under integral translations: a change
of that column affects the head of \eqref{eq:source-18} only at
depth $s_i$, and preserves the trailing integrality tests.  The
Fourier transform therefore has integral support and introduces no
nonunit normalization at $p$.  Partition the remaining position
tests by cosets and decompose them into products of unary tests.
Their coefficients retain the same supremum bound, so
Lemma~\ref{theta:cut-moments} applies to their monomial moments.

The column-degree filtration of the differential bundle identifies
these moments with the mixed coefficients.  Binary columns have
the normalized unary moments, and the hyperbolic column contributes
integral position powers.  Logarithmic frames and prime-to-$p$
Jacobi torsion evaluations split the required differential
sequences integrally.  These are the mixed-expansion identities of
\cite[Lemma 4.5]{C}; the argument
above checks the only changing input, the increasing $p$-depth.

At moving and fixed finite places choose the boundary
representatives by approximation, with the required integral
conditions and the deep $p$-condition simultaneously.  Detecting
Jacobi torsions remain prime to $p$.  Source integration has a fixed
volume denominator because the determinant equivariance cancels
the character on the fixed source compact, including at $p$.
Thus all expansion coefficients have one fixed denominator.
Saturation on the canonical branch, supplied at every depth by
Lemma~\ref{theta:canonical-levels}, gives the assertion for the
sections themselves.  This concerns the canonical ordinary lattice;
it makes no assertion about an unrelated generic-level lattice.
\end{proof}

\subsection{Toric periods and Fourier cancellation at \texorpdfstring{$p$}{p}}

We next identify the limit of the toric period entering the constant
term.  At a split place write $a(x)=\operatorname{diag}(x,1)$ and
$n(y)=\left(\begin{smallmatrix}1&y\\0&1\end{smallmatrix}\right)$.
In Mellin notation put $\chi_p^{\mathrm{Mell}}=\eta_p^{-1}$, and
choose the Kirillov function
\[
 W_{\chi_p^{\mathrm{Mell}}}(a(x))
   =(\chi_p^{\mathrm{Mell}})^{-1}(x)
      \mathbf1_{\Z_p^\times}(x).
\]
This compactly supported function occurs in every generic local
elliptic representation.  With unit-group volume one its central
Mellin integral against $\chi_p^{\mathrm{Mell}}$ is $1$, whether
that character is ramified or unramified.  Its unit equivariance
cancels the torus weight, so integration again has fixed volume.

Let $W_0(a(x))=\mathbf1_{\Z_p^\times}(x)$, the canonical
$p$-depletion test, and let $b_p$ be large enough for the character
on units.  Finite Fourier inversion in the Kirillov model gives
\begin{equation}
 W_{\chi_p^{\mathrm{Mell}}}
 =p^{-b_p}\sum_{j\bmod p^{b_p}}
   \left(\sum_{\substack{x\bmod p^{b_p}\\p\nmid x}}
    (\chi_p^{\mathrm{Mell}})^{-1}(x)
       \psi_{\mathrm{add}}(-xj/p^{b_p})\right)
       n(j/p^{b_p})W_0.
 \tag{19}\label{eq:source-19}
\end{equation}
The translations are right translations.  Their meaning in CM
coordinates is needed before one estimates this expression.

\begin{lemma}[CM translations and bounded Fourier projection]
\label{theta:fourier-measure}
Choose connected and etale Tate coordinates $e_p,f_p$ at a split
ordinary CM point.  The translation $n(j/p^b)$ gives the lattice
\[
 \Z_pe_p+\Z_p\bigl(f_p+(j/p^b)e_p\bigr).
\]
It corresponds, up to a fixed sign and a unit root-basis choice,
to the Serre--Tate parameter $\zeta_{p^b}^{cj}$ with
$c\in\Z_p^\times$, and preserves the connected differential frame
integrally.  Suppose $F(z)=\int_{\Z_p}z^y\,d\mu(y)$ for a bounded
measure $\mu$, and $\chi$ has conductor dividing $p^b$ on units.
For every $\delta\in\Z_p^\times$,
\begin{equation}\label{theta:fourier-cancellation}
 p^{-b}\sum_{j\bmod p^b}\sum_{x\in(\Z/p^b)^\times}
     \chi^{-1}(x)\zeta_{p^b}^{-xj}
       F(\zeta_{p^b}^{c\delta j})
   =\int_{\Z_p^\times}\chi^{-1}(c\delta y)\,d\mu(y).
\end{equation}
Its norm is at most $\|\mu\|$, independently of $b$.
\end{lemma}

\begin{proof}
An ordinary deformation is the extension of the etale group by the
multiplicative group classified by its Serre--Tate unit $z$.
For $z=\zeta_{p^b}^j$, pulling back on the etale generator by $p^b$
makes the extension split.  Relative to that rational splitting,
the compatible roots $\zeta_{p^{a+b}}^j$ of $z$ shift an etale
lift at level $a$ by $j/p^b$ in the connected coordinate.  This
gives the displayed lattice and parameter; changing the compatible
root basis or reciprocity sign gives the fixed unit $c$.

The lattice intersects the connected rational line in precisely
$\Z_pe_p$, and its image in the etale rational line is precisely
$\Z_pf_p$.  Thus the rational comparison transports the connected
lattice integrally, not by multiplication by a power of $p$.
One can also realize it by canonical connected quotients followed
by etale quotients with the chosen roots; dividing the resulting
isogeny by the common scalar gives the same integral connected
comparison.  Its multiplicative differential changes by a unit.
Canonical levels in $W_0$ remain canonical, and the rational CM
splitting agrees with the ordinary unit-root splitting by the CM
idempotents.  In particular no factor $p^{bm}$ appears in a weight
comparison.

Substitute the measure expression for $F$ into the left side of
\eqref{theta:fourier-cancellation}.  The finite sum over $j$ is
$p^b$ when $x\equiv c\delta y\pmod{p^b}$ and zero otherwise.
There is exactly one such unit $x$ if $y$ is a unit, and none if
$y$ is not.  Division by $p^b$ gives the right side.  This proves
the equality and the bound without estimating the Fourier
coefficients separately.
\end{proof}

Let $P_{0,i}(m)$ and $P_{1,i}(m)$ be the raw toric periods for the
forms of $E$ and $E^{D'}$, respectively.  Their definition is the
CM toric sum in \cite[Section 5]{C}, with
\eqref{eq:source-19} at $p$, the raw factors $q-1$ and translations
$n(1/q)$ at the moving primes, and the oriented fixed local tests
specified below.  Raw means that conductor-class sums are not
averaged by their growing orders.

\begin{proposition}[The toric weight limits]\label{theta:toric-limit}
At every fixed stage $i$, the periods $P_{j,i}(m)$ have coefficientwise
limits as positive $m$ tends to zero on the chosen $p$-adic branch
and tends to infinity ordinarily.  The limits, as group polynomials,
have uniformly bounded coefficients as $i$ varies.  Up to a fixed
nonzero scalar and group-like units, the limit of $P_{0,i}(m)$ is
the stage image of the undepleted series $b_w^0$.
Consequently, when $b_w^0(z_0)=0$, weights may be chosen at each
stage with arbitrary prescribed additional precision so that
\[
 P_{0,i}(m_i)(z_0)\longrightarrow0,
 \qquad \sup_i|P_{1,i}(m_i)(z_0)|_p<\infty.
\]
\end{proposition}

\begin{proof}
First keep $i$ and $b_p$ fixed.  The CM/Katz comparison and
Proposition~\ref{pair:primitive} identify the weight limit for the
first form with $F_0=d_{\mathrm{mod}}^{-1}f^{[p]}$ in integral
Igusa frames.  The CM representatives at the moving tame conductor
can be chosen with $p$-component one.  Transport reference algebraic
CM frames through their toric calculations and use the corresponding
Igusa frames for the derivative limit.  All degree-dependent
comparison factors are then fixed unit powers.  The translations
in \eqref{eq:source-19} preserve those frames by
Lemma~\ref{theta:fourier-measure}.  Thus the branch conditions are
compatible; no new nonunit raised to the weight is introduced.

For the disk indexed by $\mathfrak a$, a prime-to-$p$ isogeny
changes the root parameter to $\zeta_{p^{b_p}}^{c\delta_{\mathfrak a}j}$.
The unit action $a(u)n(j/p^{b_p})$ changes the index to $uj$.
CM reciprocity therefore expresses the unit character in
\eqref{eq:source-19} through the exponent map
$\eta_\Gamma:\Z_p^\times\to\Gamma$ of Section~\ref{sec:pair}.
Formula~\eqref{theta:fourier-cancellation} gives the inverse moment
label
$[-\eta_\Gamma(y)-\eta_\Gamma(\delta_{\mathfrak a})]$, up to the common group-like
unit coming from $c$.  The represented torus class contributes
$[t_{\mathfrak a}]\psi_s(\mathfrak a)$.  Summing the disk expressions
is exactly the formula for the undepleted $b_s^0$ in
Proposition~\ref{pair:interpolation}.  Changes of basepoint, root
basis, or non-$p$ orientation introduce only its fixed sign and
translation units.  This establishes the identification including
the global scalar labels.

For the second form the identical local Fourier calculation suffices
for a bound.  Its fixed toric tests away from $p$ are finite
expressions of translates of the primitive form with fixed
differential constants.  Its inverse-derivative functions on the
ordinary CM disks give bounded measures.  Unit restriction,
prime-to-$p$ translations, and the raw conductor sums preserve a
common bound.  The first form has the same bound from the explicit
$b_s^0$ formula.  The two limits are therefore uniformly bounded
group polynomials.

There is no assertion here that inverse Fourier transform at all
stages has uniformly bounded norm.  At fixed $i$, its denominator
$p^{b_p}$ is fixed, and weight convergence may be taken to any
precision before letting $i$ grow.  Equivalently, characterwise
convergence in the fixed characteristic-zero group algebra implies
coefficientwise convergence with a finite stage-dependent loss,
which is paid by taking the weight sufficiently close to zero.
The resulting limiting polynomials have the uniform bound proved
above.  Substitution at the interior point respects their limits.
Since $b_w^0(z_0)=0$, the stage limits for the first period tend to
zero there; the second limits remain bounded.  Approximating each
limit to any prescribed stage precision proves the assertion.
\end{proof}

\subsection{The binary comparison and its denominators}

At zero Fourier--Jacobi index the mixed expansion sets the
hyperbolic position equal to zero.  Every polynomial containing a
hyperbolic-column minor vanishes.  The remaining determinant
monomial gives a binary theta value on $\mathrm U(B)$.
At a split finite place its local intertwiner, on the bare source
space, is
\begin{equation}\label{theta:binary-map}
 J_v(\Phi,z)=\int_{\mathrm{GL}_2(F_v)}
  \Phi({}^th)|\det h|_v\,
  \pi_v^\vee(h)\eta_v^{-1}(\det h)z\,d^\times h,
 \qquad \operatorname{vol}\mathrm{GL}_2(\mathcal O_{F_v})=1.
\end{equation}
The output uses inverse transpose.  This is the central
Godement--Jacquet zeta map \cite{GodementJacquet}, with rational
continuation when required.
The compact-model multiplicity statement of
\cite[Section 5, ``The binary scalar and its multiplicity'']{C} expresses the global
binary map as one scalar times the local maps.  Its hypotheses here
are the spherical nonsplit components and cuspidality of the base
change already checked above.

For precision about the comparison being used, write
$\chi=\eta_v^{-1}$ in the distinguished split coordinate.  The
matrix zeta function and the toric Mellin function are
\[
 \begin{split}
 Z_G(\Phi,c,\chi,s)
   &=\int_{\mathrm{GL}_2(F_v)}
       \Phi(g)c(g)\chi(\det g)|\det g|^{s+1/2}\,d^\times g,\\
 M(W,\chi,s)
   &=\int_{F_v^\times}W(a(x))\chi(x)|x|^{s-1/2}\,d^\times x.
 \end{split}
\]
Their dual functional equations use the same
$\gamma(s,\pi_v\otimes\chi,\psi_{\mathrm{add}})$:
\[
 \begin{split}
 Z_G(\widehat\Phi,c^\vee,\chi^{-1},1-s)
     &=\gamma(s,\pi_v\otimes\chi,\psi_{\mathrm{add}})
                        Z_G(\Phi,c,\chi,s),\\
 M(W^w,\chi^{-1},1-s)
     &=\gamma(s,\pi_v\otimes\chi,\psi_{\mathrm{add}})
                        M(W,\chi,s).
 \end{split}
\]
The Weyl involution uses trivial central character.  These are the
normalizations in \cite[Equations (5.8) and (5.10)]{C}.  They hold for arbitrary
smooth twists, including special local representations.  In the
bilinear pairing there is one functional-equation factor, with the
same orientation on the matrix and toric sides.

The global identity is the equal-rank, anisotropic binary
Siegel--Weil/Rallis identity \cite{Yamana,GanQiuTakeda} combined
with the two Waldspurger identities \cite{Waldspurger}, with the
specific normalization of \cite[Equation (5.6)]{C}.
The source groups are
compact, their base-change parameters cuspidal and tempered, their
zero-space theta lifts vanish, and the torus characters have the
compatible opposite central characters in the two bilinear slots.
The central standard $L$-value factors as the two central Rankin
values for $E$ and $E^{D'}$ over $K$.  Thus these are value
identities, not residue or derivative identities.  Keeping partial
$L$-functions and raw exceptional local integrals gives the
comparison without dividing by a global $L$-value.  We record the
local changes required by our test point and increasing $p$-depth.

At $p$ the Mellin test above has value exactly one.  Since $p$ splits
completely in $L$, both local matrix gamma factors match the two
toric gamma factors, including their conductors.  They cancel
without any bound on their separate valuations.  The local
unramified-twist assumption in the fixed-depth calculation of
\cite[Section 5]{C} was used to choose its Mellin-one test; the test
\eqref{eq:source-19} supplies that property for our ramified
characters.  Duality exchanges the distinguished coordinates by
unit scales, so its remaining factors are character values on
units and fixed unit powers.

At a moving prime let $L_q^{\mathrm{fake}}$ be the degree-two
unramified central Euler factor formed from the uniformizer value
of $\chi_q^{\mathrm{Mell}}$, ignoring its tame inertia.  Put
$D_q^{\mathrm f}=(L_q^{\mathrm{fake}})^{-1}$.  The translated
spherical Whittaker expansion gives
\begin{equation}\label{theta:moving-clearing}
 \begin{split}
 H_q^{\mathrm f}
 &=D_q^{\mathrm f}(q-1)M(W,\chi_q^{\mathrm{Mell}},1/2)\\
 &=D_q^{\mathrm f}
      \sum_{a\in\F_q^\times}\chi_q^{\mathrm{Mell}}(a)
                    \psi_{\mathrm{add}}(a/q)
   +(1-D_q^{\mathrm f})
      \sum_{a\in\F_q^\times}\chi_q^{\mathrm{Mell}}(a).
 \end{split}
\end{equation}
Both factors have bounded group coefficients.  Our uniformizer
choice makes $D_q^{\mathrm f}(z_0)$ tend to a nonzero value.
The additive sum in \eqref{theta:moving-clearing} is a unit at the
test: modulo the maximal ideal all the group elements specialize
to one, and its augmentation is $-1$.  The plain unit sum tends
to zero at the test.  It is an integral multiple of the norm
element of its own inertia group; the geometric-sum formula gives
convergence to zero when $u_{j0}\ne0$, and at $u_{j0}=0$ its value
is an integral multiple of the group order, which also tends to
zero.  Hence
$H_q^{\mathrm f}(z_0)$ is nonzero with uniformly bounded valuation.
The additive sums may lie in varying coefficient fields; no
convergence of those unit sums in a fixed finite field is required.
The period comparison uses these full raw factors.  No inverse
of $q-1$ or conductor projector remains.

At a fixed split $q\ne p$, take the oriented toric newvector.  Its
Mellin integral is
$L(1/2,\pi_{j,q}\otimes\chi_q^{\mathrm{Mell}})$, $j=0,1$,
up to fixed volume, unit, and monomial factors with bounded inverse
at $z_0$.  Kill fixed $\mathcal C^m$ inertia by the weight branch.
Divide the binary map \eqref{theta:binary-map} at $v\mid q$ by
$L(1/2,\pi_{F,v}^\vee\otimes\chi_v^{\mathrm{Mell}})$.
For fixed Schwartz and matrix-coefficient data, the
Godement--Jacquet zeta identity says that coefficients of this
divided map are Laurent polynomials in the unramified parameter.
Their finitely many coefficients have a fixed denominator.
Local base change identifies the product over $v\mid q$ of the
binary Euler factors with the product of the two elliptic factors.
This includes
ramified quadratic base change at $q$, since $K_q$ is split.
Induction of local epsilon factors matches the gamma products,
up to the fixed lambda factor for the chosen additive characters.
Thus the Euler factors of the toric Mellins cancel those moved
into the binary scalar.  This is a rational identity before
evaluation, so it holds even when an inverse Euler polynomial
vanishes at $z_0$; no inversion of that vanishing polynomial is
made after specialization.

At fixed places where $K_q$ is nonsplit the variable scalar is
locally trivial: it is unramified and anticyclotomic of $p$-power
order, with $p$ odd.  The fixed nonzero-pairing tests and no-pole
estimates of \cite[Section 5]{C} apply, including the companion's
ramified toric functional.  At all remaining places use spherical
normalizations.

Finally, the determinant Fock norm of degree $m-1$ has ratio
$(m-1)!m!$ to the vacuum norm.  The raised quaternionic weight-two
vector has the same ratio.  These factorials cancel in the
archimedean comparison.  The CM period normalization has weight
$2m$ at each real embedding on both sides.  The remaining frame
factors are fixed algebraic unit powers, and global measures,
adjoint norms, and zeta factors contribute fixed scalars.  Their
algebraicity follows by evaluating any nonzero normalized
comparison; if no such comparison exists, the cross-multiplied
numerator identity is already zero.  This is the algebraic,
bilinear comparison, so the squares below are algebraic squares.

\begin{proposition}[Vanishing of all constant terms]
\label{theta:boundary-small}
Assume $b_w^0(z_0)=0$, and choose the auxiliary tests and local
parameters of Lemma~\ref{theta:test-point}.  After the fixed clearing
of Proposition~\ref{theta:integral-family}, weights $m_i$ may be
chosen so that all constant Fourier--Jacobi coefficients of
$\Theta_{i,m_i}$ at $z_0$ tend uniformly to zero.  The weights tend
to zero on the fixed $p$-adic branch and to infinity ordinarily.
At each fixed stage they can satisfy arbitrarily strong additional
weight-approximation requirements.  The assertion holds
simultaneously for the finitely many fixed variants used for
nonconstant content and for their integral prime-to-$p$ translates.
\end{proposition}

\begin{proof}
Let $a_{i,m}$ be any constant coefficient in an integral ordinary
frame.  It is the binary scalar times the local intertwiners
applied to the corresponding projected tests.  Square this
expression and use the bilinear comparison just described.
After cross-multiplication it has the form
\begin{equation}\label{theta:cleared-boundary}
 A_{i,m}\,a_{i,m}^{2}
     =B_{i,m}\,P_{0,i}(m)^2P_{1,i}(m)^2.
\end{equation}
Here $A_{i,m}$ and $B_{i,m}$ have uniformly bounded group
coefficients, and $A_{i,m}(z_0)$ is nonzero with valuation bounded
above, when $m$ is sufficiently close to zero at stage $i$.
We verify that this bound is uniform in the cusp, rather than
only for the reference binary test.

At a moving prime the projected test has integral support,
depends on reduction modulo $q$, and is determinant-equivariant
under the full source compact.  On a deficient-rank stratum the
stabilizer determinant surjects onto the inertia group.  Its group
coefficient is therefore an integral multiple of the norm element:
all coefficients are equal, without division by the group order.
The remaining untwisted-unit integral is a sum over image
lattices.  Express an exact reduced line or zero image by
subtraction from containment conditions.  Each containment sum is
a translate and a power of $q$ times the full-lattice fake Euler
series.  Thus $D_q^{\mathrm f}$ clears its denominator, while
$H_q^{\mathrm f}$ clears the raw toric factors.  Their bounded
powers have the bounded evaluated valuations just proved.  This
argument also applies to bounded integral sums of pivot
translates; powers of $q$ are $p$-units.

At $p$ the canonical projection has integral matrix support, by
the Fourier-support calculation in
Proposition~\ref{theta:integral-family}.  Its determinant
equivariance cancels the source twist on a fixed compact under
which the bare vector is invariant.  Expanding
\eqref{theta:binary-map} by determinant valuation gives the factor
\[
                 p^{(m-1)v_p(\det h)}.
\]
Right-coset volume denominators are fixed, even though the target
depth grows.  The translates of the fixed compact automorphic
function have bounded sup norm.  For sufficiently large positive
$m$ this integral-matrix series converges with a uniform bound.
Its finite determinant truncations are the
Godement--Jacquet series and hence compute the same rational zeta
operator.  No projector onto a deep $p$-character space is used.

At a fixed split place the divided intertwiners above have Laurent
polynomial bounds on the finite spaces of tests.  Integral cusp
representatives and their fixed variants give the same bound.
Fixed nonsplit denominators have bounded inverse at $z_0$;
unmarked places have spherical normalization.  Boundary
approximation chooses representatives integral at the designated
finite places, and at $p$ preserving the deep canonical flag.
Compact-model evaluations and frame changes therefore introduce
only the fixed factors and unit character values already
accounted for.  These observations prove the claimed uniformity
of \eqref{theta:cleared-boundary}.

The period identities first hold at actual finite characters.
Those characters separate the characteristic-zero group algebra,
so the cross-multiplied identity holds coefficientwise.  The
uniform bounds permit substitution at $z_0$ to the increasing
precision for which the finite-group relations hold.  By
Proposition~\ref{theta:toric-limit}, its right side tends to zero,
and the valuation of the nonzero left multiplier stays bounded.
Thus $v_p(a_{i,m_i}(z_0))\to\infty$, uniformly in the coefficient.
Taking a square root of this valuation inequality merely divides
the required precision by two.

For clarity, the parameters are chosen in the following order.
Fix the finite group stage and its required $p$-depth.  At that
stage take the positive weight sufficiently close to zero to pay
every finite Fourier-inversion loss and approximate the two
toric limits to the desired precision.  Increase its ordinary
size to meet the convergent-series and holomorphy requirements.
Then let the group orders increase and retain a precision
tending to infinity more slowly than their relation precision
and the resulting constant-term precision.  Fixed losses are
absorbed in this retained precision.  Additional finitely many
weight congruences and lower bounds can be imposed at the first
step.  In particular the degree choices used to retain
nonconstant content in Section~\ref{sec:extraction} are compatible
with this construction.
\end{proof}
